\documentclass[a4paper,10pt]{article}
\usepackage[francais]{babel}
\usepackage{url}

\selectlanguage{francais}
\input{letter-ensiie.sty}

\selectlanguage{francais}
\enCouleur

\begin{document}


\name{St\'ephane \textsc{Genaud}}
\fonction{Professor at ENSIIE}
\telephoneillkirch{+33 (0)3688 54542}
\email{genaud@ensiie.fr}

\begin{letter}[Recommendation de Luke Bertot]{à}%
{Ecole Doctorale MSII}


Chers Collègues,

Je  connais Luke  Bertot  depuis sa  première année  de  scolarité dans  l'école
d'ingénieur  en  informatique   ENSIIE,  dont  je  dirige   l'antenne  locale  à
Strasbourg, et dans laquelle j'interviens également comme enseignant. Je suis
également enseignant-chercheur et travaille avec Luke sur un projet ANR depuis
janvier 2014. Je donne des éléments d'appréciation de la candidature de Luke
à ces deux titres.

\textbf{Résultats Scolaires}

La formation d'ingénieur de l'ENSIIE est dispensée à la fois sur le site
historique principal, à Evry, et à son antenne de Strasbourg selon les modalités
suivantes: un tronc commun se déroule sur trois semestres, avec des contenus 
pédagogiques et des examens communs aux deux sites, puis les semestres 4 et 5
sont constitués à 80\% d'options axées sur les spécialités recherche des sites.
Le dernier semestre est un stage long de type professionnel ou recherche. 

En tant  que directeur  de l'école,  j'ai obtenu des  services de  scolarité les
résultats scolaires  de l'ensemble des élèves  sur les deux sites.  Sur le tronc
commun (S1/S2/S3), la moyenne générale (pondérée par les ECTS) de Luke permet de
comparer   ses  résultats   à  l'ensemble   de  la   promotion  (106   puis  124
élèves-ingénieurs).

Sur le parcours  d'option (S4/S5) la comparaison est plus  difficile étant donné
que même  sur un site,  les options choisies par  les élèves peuvent  varier. Le
tableau fait apparaître cependant un résultat qui donne la tendance de ses
résultats scolaires.

\begin{center}
\begin{tabular}{ccc}
semestre & rang/effectif & moyenne générale\\
\hline \\
S1 & 12/106 & 14,7/20\\
S2 & 20/106 & 15,1/20\\
S3 & 31/124 & 13,7/20\\
S4 & 7/17   & 12,1/20\\
S5 & 4/17   & 14,8/20\\
\hline \\
\end{tabular}
\end{center}

Globalement, Luke se classe dans le premier quart d'une promotion recrutée sur
un concours d'entrée sélectif (concours Centrale-Supélec).


\textbf{Motivation pour la Recherche}

J'ai recruté  Luke Bertot comme  ingénieur de  recherche pour travailler  sur un
projet ANR dont je suis porteur local. Luke participe pleinement à l'activité de
l'équipe, en  collaborant avec les  doctorants et  permanents, et en  ayant déjà
participé à deux réunions nationales de l'ensemble du projet.

Cette  période   de  six  mois  durant   laquelle  Luke  a  travaillé   sur  nos
problématiques de recherche m'a pleinement conforté dans l'idée que Luke possède
toutes  les  qualités  pour mener  à  bien  un  travail  de thèse.   Il  possède
indiscutablement de solides connaissances  scientifiques générales, la curiosité
indispensable à  ce métier, ainsi  qu'une formation initiale en  informatique en
adéquation  avec les  tâches qu'il  aura  à mener  pour valider  son travail  de
recherche.


Pour ces raisons, je suis convaincu que Luke Bertot est un excellent candidat
pour l'obtention d'un contrat doctoral.




\end{letter}
\begin{flushright}
%\includegraphics[width=.2\textwidth]{signgenaud.jpg}
\end{flushright}
\end{document}
